\documentclass[10pt,a4paper]{amsart}
\usepackage[margin=19mm]{geometry}
\usepackage{amsmath,amssymb,mathtools}
\usepackage[scr=rsfs,cal=euler]{mathalfa}
\usepackage{xfrac}
\usepackage[unicode,hidelinks,bookmarks=false]{hyperref}
\newcommand{\ie}{i.e.\ }
\newcommand{\cf}{cf.\ }
\newcommand{\resp}{respectively\ }
\newcommand{\Subsection}{\subsection}
\providecommand{\nobreakdash}{\nobreak\hbox{-}}
\setlength{\emergencystretch}{2em}
\multlinegap0pt
\usepackage{bm}
\usepackage{bbm}
\usepackage{dsfont}
\usepackage{xspace}
\usepackage{cancel}
\usepackage{mathtools}
\usepackage{amsthm}
\makeatletter
\@ifundefined{theorem}{\newtheorem{theorem}{Theorem}}{}
\@ifundefined{lemma}{\newtheorem{lemma}[theorem]{Lemma}}{}
\makeatother
\newcommand\E{{\mathbb{E}}}  
\newcommand\N{{\mathbb{N}}}
\renewcommand\P{{\mathbb{P}}}
\title[On the First Non-Universal Term in Random Polynomial Real Ze]{On the First Non-Universal Term in Random Polynomial Real Zeros}
\author{Phuc Lam, Oanh Nguyen}
\date{Source: arXiv:2509.12170v1 (CC BY 4.0)}
\begin{document}
\maketitle

\noindent\emph{Source and licence.} On the First Non-Universal Term in Random Polynomial Real Zeros. Version arXiv:2509.12170v1 (CC BY 4.0), distributed under the Creative Commons Attribution 4.0 International licence, as recorded in the arXiv licence metadata for this exact version and shown on the abstract page https://arxiv.org/abs/2509.12170v1. Full rights notice: licenses/arxiv-2509.12170-CC-BY-4.0.txt.\par\smallskip

\noindent\textit{Editorial scope.} Counted unit: the Lemma whose source label is \texttt{lem:ConvFarFrom01Cont} in the article (source0.tex lines 694--696, source label \texttt{lem:ConvFarFrom01Cont}), delivered together with the article's own complete original proof of it (source0.tex lines 698--734, one \texttt{proof} environment of 37 lines). No mathematics is written, repaired or re-derived: statement and proof are the article's own text line for line. Required context retained uncounted: lem:ConvFarFrom01 (lines 354--357); lem:momentsOfNumRoots (lines 493--496); lem:SkorokhodConstruction (lines 651--657). Every definition, display and label of those lines is retained unchanged. Omitted from this package and not redistributed: 1--353, 358--492, 497--650, 658--693, 697--697, 735--802. Nothing inside a retained excerpt is omitted or altered: every retained body is delivered exactly as the article's lines are written, so no omission receipt of any kind accompanies this package. Citations: the retained excerpts carry no citation command at all, so no bibliography entry of the article is transcribed and no reference is redistributed. Reference adaptation: none was needed and none was made; every reference inside the delivered unit resolves inside it, so no unresolved reference or citation is produced. Printed slips kept literal: no printed word, symbol, hypothesis or number of the counted statement or of its proof was altered. Layout and class: the article's own class is \texttt{amsart} with packages including amsmath, amsthm, amsfonts, amssymb, verbatim, graphicx, mathtools, geometry, url, enumerate, hyperref, ulem, color; the wrapper is the lane's pinned \texttt{amsart} editorial wrapper at 10pt on A4, which restates the theorem-like environments the retained text needs and adds the article's own macro definitions that the retained text needs (\texttt{\textbackslash E}, \texttt{\textbackslash N}, \texttt{\textbackslash P}), with extra wrapper packages (bm, bbm, dsfont, xspace, cancel, mathtools). The delivered numbering is the unit's own. Assets: no retained span contains a figure environment, a graphics-inclusion command, a PDF-page inclusion, a table or a TikZ picture; no image file is copied into this package, the package declares no asset dependency and no image file is redistributed with it. Licence: version arXiv:2509.12170v1 of this article is distributed under the Creative Commons Attribution 4.0 International licence, as recorded in the arXiv licence metadata for this exact version and shown on the abstract page https://arxiv.org/abs/2509.12170v1. A full rights notice is delivered at licenses/arxiv-2509.12170-CC-BY-4.0.txt. Characters: every retained excerpt is pure ASCII, so the delivered engine, XeLaTeX with its default Latin Modern text font, reports no missing character.
\begin{lemma}\label{lem:ConvFarFrom01}
    Fix any fixed $r<R$ in $ (0, 1)$, we have
    $$\lim_{n \rightarrow \infty} \E N_{P_n}([r, R]) = \E N_{P_{\infty}}([r, R]) < \infty.$$
\end{lemma}

\begin{lemma}\label{lem:momentsOfNumRoots}
	Denote $N_n := N_{P_n}([-R, R])$. Then for all $\kappa \ge 1$, there exists a constant $M = M(R, \kappa, \varepsilon_0, M_0) \in (0,\infty)$ such that for all $n$, $$ \E N_n^{\kappa} \le M,$$
	where $\varepsilon_0, M_0$ are such that $\E|\xi_k|^{2 + \varepsilon_0} \le M_0$ for all $k$. 
\end{lemma}

\begin{lemma}\label{lem:SkorokhodConstruction}
	Assume that the sequence of probability measures $\mu_{n}$ converges to $\mu_\infty$ in distribution, then we can construct $(\xi_{k, n})_{k, n\in \N}$ and $(\xi_{k, \infty})_{k\in \N}$ such that the following holds.
	\begin{enumerate}
		\item For each $n\in \N\cup \{\infty\}$, the random variables $\xi_{0, n}, \xi_{1, n}, \dots$ are i.i.d. with distribution $\mu_{n}$;
		\item For each $k$, $\xi_{k, n}$ converges to $\xi_k$ almost surely as $n\to\infty$.
	\end{enumerate}
\end{lemma}

\begin{lemma}\label{lem:ConvFarFrom01Cont}
For all $0 < r < R < 1$, we have $$\lim_{n \rightarrow \infty}\E N_{Q_{n}}[r, R] = \E N_{Q_\infty}[r, R].$$
\end{lemma}

\begin{proof}
	Firstly, we show that $Q_n$ convergence to $Q_\infty$ in finite-dimensional distributions. In particular, we show that for all $l \in \mathbb{N}$ and $z_1, \dots, z_l \in D_{R}$, as $n \rightarrow \infty$,
	\begin{equation*}
		\begin{pmatrix}
			\text{Re } Q_n(z_1) \\
			\text{Im }  Q_n(z_1) \\
			\vdots \\
			\text{Re } Q_n(z_l) \\
			\text{Im } Q_n(z_l)
		\end{pmatrix} \xrightarrow[]{d} \begin{pmatrix}
			\text{Re } Q_\infty(z_1) \\
			\text{Im } Q_\infty(z_1) \\
			\vdots \\
			\text{Re } Q_\infty(z_l) \\
			\text{Im } Q_\infty(z_l)
		\end{pmatrix}.
	\end{equation*}
	To this end, we construct the coefficients of $Q_n$ and $Q_\infty$ according to Lemma \ref{lem:SkorokhodConstruction}. With this construction, it suffices to show that for all $z \in D_R$, a.s. $Q_n(z) \to Q_\infty(z)$.
	
	Consider the event $\mathcal  A = \left\{\sum_{k \ge 0} \left(\sup_n |\xi_{k, n}|\right) z^k < \infty \right\}$. On $\mathcal A$, we have $Q_n (z) \to Q_\infty(z)$ by Dominated Convergence Theorem, so it remains to show that $\P(\mathcal A) = 1$. If $U_k$ is uniformly distributed on $[0, 1]$, then for all $\varepsilon > 0$, 
	\begin{align*}
		\P\left( \sup_n |\xi_{k, n}| \ge (1 + \varepsilon)^k  \right) &= \P\left( \exists n: |\xi_{k, n}| \ge (1 + \varepsilon)^k \right) \\
		&\le\P\left( \exists n: \xi_{k, n} \ge (1 + \varepsilon)^k \right) + \P\left( \exists n: \xi_{k, n} \le - (1 + \varepsilon)^k \right).
\end{align*}
If $n$ is such that $\xi_{k, n} := F_n^{-1}(U_k) \ge (1 + \varepsilon)^k$, then $U_k \ge F_n((1 + \varepsilon)^k) \ge 1 - (1 + \varepsilon)^{-2k}$, where the last equality follows from Chebyshev's inequality and the fact that $F_n$ has unit variance. Thus,
\begin{align*}
		\P\left( \exists n: \xi_{k, n} \ge (1 + \varepsilon)^k \right) &\le \P\left( U_k \ge  1 - (1 + \varepsilon)^{-2k}  \right) \le (1 + \varepsilon)^{-2k}. 
\end{align*}
We obtain the same bound for $\P\left( \exists n: \xi_{k, n} \le - (1 + \varepsilon)^k \right)$. Thus,
\begin{align*}
		\P\left( \sup_n |\xi_{k, n}| \ge (1 + \varepsilon)^k  \right) &\le 2( 1+ \varepsilon)^{-2k},  \\ 
		\sum_{k \ge 0 }\P\left( \sup_n |\xi_{k, n}| \ge (1 + \varepsilon)^k  \right) &\le 2 \sum_{k \ge 0}( 1+ \varepsilon)^{-2k} < \infty.
	\end{align*}
	By Borel-Cantelli lemma, a.s. there are finitely many $k$ such that $\sup_n |\xi_{k, n}| \ge (1 + \varepsilon)^k$. Pick $\varepsilon > 0$ sufficiently small such that $R (1 + \varepsilon) < 1$, then a.s. $\mathcal A$ happens.
	
	The remaining steps are identical to the proof of Lemma \ref{lem:ConvFarFrom01} as tightness is easily verified, and uniform integrability follows directly since all $\mu_n, \mu_\infty$ share the same parameters $\varepsilon_0, M_0$ in Lemma \ref{lem:momentsOfNumRoots}.
\end{proof}

\end{document}
